\documentclass[10pt,a4paper]{amsart}
\usepackage[margin=19mm]{geometry}
\usepackage{amsmath,amssymb,mathtools}
\usepackage[scr=rsfs,cal=euler]{mathalfa}
\usepackage{xfrac}
\usepackage[unicode,hidelinks,bookmarks=false]{hyperref}
\newcommand{\ie}{i.e.\ }
\newcommand{\cf}{cf.\ }
\newcommand{\resp}{respectively\ }
\newcommand{\Subsection}{\subsection}
\providecommand{\nobreakdash}{\nobreak\hbox{-}}
\setlength{\emergencystretch}{2em}
\multlinegap0pt
\usepackage{enumerate}
\usepackage{amssymb}
\usepackage{tikz}
\newtheorem{corollary}{Corollary}
\newtheorem{proposition}{Proposition}
\title{On Similarity Structure Groups and their W$^*$ and C$^*$-algebras}
\author{Eli Bashwinger, Patrick DeBonis}
\date{Source: arXiv:2507.18821v2 (CC BY 4.0)}
\begin{document}
\maketitle

\noindent\emph{Source and licence.} On Similarity Structure Groups and their W$^*$ and C$^*$-algebras. Version arXiv:2507.18821v2 (CC BY 4.0), distributed by arXiv under the Creative Commons Attribution 4.0 International licence, http://creativecommons.org/licenses/by/4.0/, as recorded in the arXiv licence metadata for this exact version and shown on the abstract page https://arxiv.org/abs/2507.18821v2. Full licence notice: \texttt{licenses/arxiv-2507.18821-CC-BY-4.0.txt}.\par\smallskip

\noindent\textit{Editorial scope.} Counted unit: the article's proposition prop:extension (source0.tex lines 423--425). The article's complete original proof of it is delivered with it (source0.tex lines 427--527, one \texttt{proof} environment of 101 lines). No mathematics is written, repaired or re-derived: the statement and the proof are the article's own lines, in the article's own notation, and no retained line is substituted or re-flowed.  Retained uncounted as the source-local context the proof needs: lines 403--405.  Nothing was omitted from any retained span, so the package carries no omission record.  No retained line contains a citation, so the wrapper carries no bibliography and no bibliography entry.  No retained line contains a figure environment, an image-inclusion command or drawing code, so no image is delivered and the package declares no asset dependency.  Editorial formatting only, in addition: an AMS \texttt{amsart} preamble that restates the article's own declarations and macros used by the retained lines, namely the environments \texttt{corollary}, \texttt{proposition} on separate counters, together with the standard packages those lines need.  The retained environments print in this document as Corollary 1, Proposition 1 in order of appearance, whereas the article prints the same items with the numbering of its own full document.  The complete native source of the article as distributed in the arXiv e-print for this exact version is bundled unchanged as \texttt{sources/182174/source0.tex}.
\begin{corollary}\label{cor:extension}
Let $\Gamma \leq \text{LS}(X)$ be a CSS$^*$ group. Let $B$ be an arbitrary ball in $X$ and $D$ an infinite ball in $X$ disjoint from $B$. Then there exists a ball $D_0 \subseteq D$  and an involution $h \in \Gamma $ such that $h(B) = D_0$ and $h(D_0) = B$. 
\end{corollary}

\begin{proposition}[Extension Property]\label{prop:extension}
    Let $\Gamma \leq \text{LS}(X)$ be a CSS$^*$ group and let $S_1 =\{B_1, \dots , B_n\}$ and $S_2 = \{D_1, \dots , D_n\}$ be two sets of infinite balls in $X$ that are respectively pairwise disjoint. Assume that neither $S_1$, $S_1$ or $S_1 \cup S_2$, is a complete partition of $X$. Suppose that there exists $h_i \in \text{Sim}(B_i,D_i)$ for all $1 \leq i \leq n$. Then there exists an element $g$ in $\Gamma$ such that $g\vert_{B_i} = h_i$ for all $1 \leq i \leq n$.
\end{proposition}

\begin{proof}

    If $B_i \cap D_j = \emptyset$, for all $i,j$ then we can simply take a partition $\mathcal P$ of $X$ such that $S_1, S_2 \subset \mathcal P$ and use the pasting lemma to define $g$ of $\Gamma$ satisfying the proposition to be
    
    \[ g(x) = \begin{cases} h_i(x) & x \in B_i \\
    h_i^{-1}(x) & x \in D_i \\
    x & x \in X \backslash \{S_1, S_2\}
    \end{cases}.
    \]  

    Note that since $X \backslash \{S_1, S_2\}$ is still a compact ultrametric, there exists a partition of $X \backslash \{S_1, S_2\}$ into closed balls where we are simply taking the identity similarity as is done in Corollary \ref{cor:extension}. This property is used repeatedly throughout, and we do not write down partitions where the identity similarity is being used.
    

    When there exists $i,j$ for which $B_i \cap D_j \neq \emptyset$, the general process to construct a $g \in \Gamma$ such that $g\vert_{B_i} =h_i$, is as follows. Let $B_{i_1} \cap D_{j_1} \neq \emptyset$, $i_1,j_1 \in \{ 1, \dots , n\}$, be the first non-empty intersection. If $i_1=j_1$, proceed as in Case 1 below, otherwise assume $i_1 \neq j_1$. If $B_{i_1} \subset D_{j_1}$, choose a partition $\{B_{i_1}, \tilde B_1\dots ,\tilde B_r \}$ of $D_{j_1}$ and choose a ball $A_1$ such that $A_1$ is disjoint from all $B_i$ in $S_1$ and all $D_i$ in $S_2$. Since $\Gamma$ is a CSS$^*$ group, there exists $\tilde A_1 \subseteq A_1$ such that there is a nontrivial $f_1 \in \text{Sim}(D_{j_1},\tilde A_1)$. Then, using the pasting lemma, construct an element $g \in \Gamma$ such that $g\vert_{B_i} =h_{i}$ for all $1 \leq i \leq n$, $g\vert_{\tilde B_1} = f\vert_{\tilde B_i}$ for all $1 \leq i \leq r$, $g\vert_{D_{j_1}} = f_1\vert_{B_{i_1}}h^{-1}_{i_1}$, and $g\vert_{\tilde A_1} =h_{j_1}^{-1}f_{j_i}^{-1}$. If there are no other intersections, stop and extend $g$ to be an element on $\Gamma$, by letting $g\vert_{D_i} = h_i^{-1}$ for all $i \neq i_1,j_1$ and $g$ be the identity on $X \setminus\{S_1, S_2, \tilde A\}$. Otherwise, proceed to $B_{i_2}\cap D_{j_2} \neq \emptyset$ and repeat the procedure outlined above. Note that since multiple $B_i$ can intersect $D_{j_1}$ and multiple  $D_j$ can intersect $B_{i_1}$ we explicitly write out the specific cases that will contribute to building $g$ for an arbitrary amount and types of intersections between $S_1$ and $S_2$. We also note that there are only $n$ types of intersections possible in the construction of $g$.
    
    
   
    
    By considering the case where $S_1 = \{B_1, B_2\}$ and $S_2 = \{D_1, D_2\}$, we can exhibit all of the different ways to construct $g$ depending on how $S_1$ and $S_2$ intersect. Thus, we consider the following cases. 
    
    
   


    \noindent\textbf{Case 1}: Suppose that $B_1 \subseteq D_1$ and $B_2 \cap D_j = \emptyset$ for all $j$.  If $B_1 = D_1$, this is is still covered by the disjoint setting with $h_i \in \text{Sim}(B_1,D_1)$ being the identity, so suppose $B_1 \subset D_1$ is a proper inclusion. Let $\{B_1, \tilde B_1, \dots \tilde B_r \}$ be a partition of $D_1$ into balls that includes $B_1$ and pick a ball $A$ disjoint from both $S_1$ and $S_2$. Since $\Gamma$ is a CSS$^*$ group, there exists a subball $\tilde A \subseteq A$ such that $\text{Sim}(D_1,\tilde A)$ is non-empty. Pick $f$ in $\text{Sim}(D_1,\tilde A)$. By the restriction property, we have that $f\vert_{B_1} \in \text{Sim}(B_1, f(B_1))$ and $f\vert_{\tilde B_i} \in \text{Sim}(\tilde B_i, f(\tilde B_i))$ for $1 \leq i \leq r$. Recall that by assumption we have $h_i \in \text{Sim}(B_i,D_i)$ for $i \in \{1,2\}$. By the composition property, we have that $f\vert_{B_1}h_1^{-1}f^{-1} \in \text{Sim}(\tilde A, f(B_1))$. Therefore, applying the pasting lemma again, we can define an element $g$ of $\Gamma$ satisfying the conclusion of the proposition as follows:


     \[ g(x) = \begin{cases} h_i(x) & x \in B_i, i\in\{1,2\} \\
     f\vert_{\tilde B_i}(x) & x \in \tilde B_i, 1 \leq i \leq r\\
    h_2^{-1}(x) & x \in D_2, \\
    f\vert_{B_1}h_1^{-1}f^{-1} (x) & x\in \tilde A\\
    x & x \in X \backslash \{S_1, S_2, \tilde A\}
    \end{cases}
    \]

     \noindent\textbf{Case 2}: Suppose that $B_1=D_2$, $B_2 \cap D_j = \emptyset$ for all $j$, and $D_1 \cap B_i = \emptyset$ for all $i$, then by the pasting lemma, an element $g$ of $\Gamma$ satisfying the conclusion of the proposition is

    \[ g = \begin{cases} h_i(x) & x \in B_i, i\in\{1,2\}  \\
    h_2^{-1}h_1^{-1}(x) & x \in D_1 \\
    x & x \in X \backslash \{S_1, S_2\}
    \end{cases}
    \]

    \noindent\textbf{Case 3:} 
    Suppose that $B_1, B_2 \subset D_1$, are proper inclusions. Take $\{B_1, B_2, \tilde B_1, \dots \tilde B_r \}$ to be a partition of $D_1$ into balls that includes $B_1$ and $B_2$. Take a ball $A$ such that $A \cap S_2 =\emptyset$ and $\tilde A \subseteq A$, the subball with $f \in$ $\text{Sim}(D_1, \tilde A)$. Following the same procedure as in the previous cases, we have that $f\vert_{B_1} \in \text{Sim}(B_1, f(B_1))$, $f\vert_{B_2} \in \text{Sim}(B_2, f(B_2))$, and $f\vert_{\tilde B_i} \in \text{Sim}(\tilde B_i, f(\tilde B_i))$ for $1 \leq i \leq r$. We also have $f\vert_{B_2}h_2^{-1} \in \text{Sim}(D_2,f(B_2))$ and $f\vert_{B_1}h_1^{-1}f^{-1} \in \text{Sim}(\tilde A, f(B_1))$. Now applying the pasting lemma produces the following element of $g$ which satisfies the conclusion of the proposition.  

    
     \[ g(x) = \begin{cases} h_i(x) & x \in B_i, i\in\{1,2\}  \\
     f\vert_{\tilde B_i}(x) & x \in \tilde B_i, 1\leq i \leq r\\
    f\vert_{B_2}h_2^{-1}(x) & x \in D_2, \\
    f\vert_{B_1}h_1^{-1}f^{-1} (x) & x\in \tilde A\\
    x & x \in X \backslash \{S_1, S_2, \tilde A\}
    \end{cases}
    \]



    \noindent\textbf{Case 4:}
    Suppose now that $B_1 \subset D_1$ and $B_2 \subset D_2$ are proper inclusions. Take $\{B_1, \tilde B_1, \dots \tilde B_r \}$ to be a partition of $D_1$ into balls that includes $B_1$ and $\{B_2, \hat B_1, \dots \hat B_s \}$ be a partition of $D_2$ into balls that includes $B_2$. Again, pick a ball $A$ disjoint from both $S_1$ and $S_2$ and another ball $E$ such that $E \cap \{S_1, S_2, A\} = \emptyset$. Since $\Gamma$ is a CSS$^*$ group, there exists $\tilde A \subseteq A$ and $\tilde E \subseteq E$ such that $f_1 \in \text{Sim}(D_1, \tilde A)$ and $f_2 \in \text{Sim}(D_2, \tilde E)$. Following the same procedure as in the previous cases, we have that $f_1\vert_{B_1} \in \text{Sim}(B_1, f_1(B_1))$,  $f\vert_{\tilde B_i} \in \text{Sim}(\tilde B_i, f_1(\tilde B_i))$ for $1 \leq i \leq r$, $f_2\vert_{B_2} \in \text{Sim}(B_2, f_2(B_2))$, and $f_2\vert_{\hat B_i} \in \text{Sim}(\hat B_i, f_2(\hat B_i))$ for $1 \leq i \leq s$.  We also have $f_1\vert_{B_1}h_1^{-1}f_1^{-1} \in \text{Sim}(\tilde A,f_1(B_1))$ and $f_2\vert_{B_2}h_2^{-1}f_2^{-1} \in \text{Sim}(\tilde E, f_2(B_2))$. Finally, applying the pasting lemma produces the following element of $g$ which satisfies the conclusion of the proposition.  

     \[ g(x) = \begin{cases} h_i(x) & x \in B_i, i\in\{1,2\}  \\
     f_1\vert_{\tilde B_i}(x) & x \in \tilde B_i, 1 \leq i \leq r\\
      f_2\vert_{\hat B_i}(x) & x \in \hat B_i, 1 \leq i \leq s\\
    f_1\vert_{B_1}h_1^{-1}f_1^{-1} (x) & x\in \tilde A\\
    f_2\vert_{B_2}h_2^{-1}f_2^{-1} (x) & x\in \tilde E\\
    x & x \in X \backslash \{S_1, S_2, \tilde A, \tilde E\}
    \end{cases}
    \]
    

    
    \noindent\textbf{Case 5:}
    Suppose now that $B_1 \subset D_2$ and $B_2 \subset D_1$ are proper inclusions.  Take $\{B_2, \tilde B_1, \dots \tilde B_r \}$ to be a partition of $D_1$ into balls that includes $B_2$ and $\{B_1, \hat B_1, \dots \hat B_s \}$ be a partition of $D_2$ into balls that includes $B_1$. Take a ball $A$ disjoint from both $S_1$ and $S_2$ and another ball $E$ such that $E \cap \{S_1, S_2, A\} = \emptyset$. Since $\Gamma$ is a CSS$^*$ group, there exists $\tilde A \subseteq A$ and $\tilde E \subseteq E$ such that $f_1 \in \text{Sim}(D_1, \tilde A)$ and $f_2 \in \text{Sim}(D_2, \tilde E)$. Following the procedure as in the previous cases, we have that $f_1\vert_{B_2} \in \text{Sim}(B_2, f_1(B_2))$,  $f_1\vert_{\tilde B_i} \in \text{Sim}(\tilde B_i, f_1(\tilde B_i))$ for $1 \leq i \leq r$, $f_2\vert_{B_1} \in \text{Sim}(B_1, f_2(B_1))$, and $f_2\vert_{\hat B_i} \in \text{Sim}(\hat B_i, f_2(\hat B_i))$ for $1 \leq i \leq s$.  We also have $f_2\vert_{B_1}h_1^{-1}f_1^{-1} \in \text{Sim}(\tilde A,f_2(B_1))$ and $f_1\vert_{B_2}h_2^{-1}f_2^{-1} \in \text{Sim}(\tilde E, f_1(B_2))$. Finally, applying the pasting lemma produces the following element of $g$ which satisfies the conclusion of the proposition.  

    \[ g(x) = \begin{cases} h_i(x) & x \in B_i, i\in\{1,2\}  \\
     f_1\vert_{\tilde B_i}(x) & x \in \tilde B_i, 1 \leq i \leq r\\
      f_2\vert_{\hat B_i}(x) & x \in \hat B_i, 1 \leq i \leq s\\
    f_2\vert_{B_1}h_1^{-1}f_1^{-1} (x) & x\in \tilde A\\
    f_1\vert_{B_2}h_2^{-1}f_2^{-1} (x) & x\in \tilde E\\
    x & x \in X \backslash \{S_1, S_2, \tilde A, \tilde E\}
    \end{cases}
    \]


    \noindent\textbf{Case 6:}
    Suppose now that $B_1 \subset D_1$ and $D_2 \subset B_2$ are proper inclusions and take $\{B_1, \tilde B_1, \dots \tilde B_r \}$ to be a partition of $D_1$ into balls that includes $B_1$ and $\{D_2, \tilde D_1, \dots \tilde D_s \}$ be a partition of $B_2$ into balls that includes $D_2$. Again, pick a ball $A$ disjoint from both $S_1$ and $S_2$ and another ball $E$ such that $E \cap \{S_1, S_2, A\} = \emptyset$. Since $\Gamma$ is a CSS$^*$ group, there exists $\tilde A \subseteq A$ and $\tilde E \subseteq E$ such that $f_1 \in \text{Sim}(D_1, \tilde A)$ and $f_2 \in \text{Sim}(B_2, \tilde E)$. Following the same procedure as in the previous cases, we have that $f_1\vert_{B_1} \in \text{Sim}(B_1, f_1(B_1))$,  $f_1\vert_{\tilde B_i} \in \text{Sim}(\tilde B_i, f_1(\tilde B_i))$ for $1 \leq i \leq r$, $f_2\vert_{D_2} \in \text{Sim}(D_2, f_2(D_2))$, and $f_2\vert_{\tilde D_i} \in \text{Sim}(\tilde D_i, f_2(\tilde D_i))$ for $1 \leq i \leq s$.  We also have that $f_1\vert_{B_1}h_1^{-1}f_1^{-1} \in \text{Sim}(\tilde A,f_1(B_1))$ and $f_2\vert_{D_2}h_2^{-1}f_2^{-1} \in \text{Sim}(\tilde E, f_2(D_2))$. Finally, applying the pasting lemma produces the following element of $g$ which satisfies the conclusion of the proposition.  

     \[ g(x) = \begin{cases} h_i(x) & x \in B_i, i\in\{1,2\}  \\
     f_1\vert_{\tilde B_i}(x) & x \in \tilde B_i, 1 \leq i \leq r\\
        f^{-1}_2\vert_{\tilde D_i}(x) & x \in f_2(\tilde D_i), 1 \leq i \leq s\\
    f_1\vert_{B_1}h_1^{-1}f_1^{-1} (x) & x\in \tilde A\\
    f_2h_2^{-1}f_2^{-1}\vert_{D_2} (x) & x\in \tilde E\\
    x & x \in X \backslash \{S_1, S_2, \tilde A, \tilde E\}
    \end{cases}
    \]

    \noindent \textbf{Case 7:} Suppose now that $B_1 \subset D_2$ and $D_1 \subset B_2$ are proper inclusions. Then an element $g \in \Gamma$ that satisfies the proposition can be built by following a straightforward combination of Case 5 and Case 6. 
    
    In the general case of $S_1 \cap S_2 \neq \emptyset$ any number of $B_i's$ can be contained in any number of $D_j's$ and vice versa. However, following the initial general procedure, we can always construct an element $g \in \Gamma$ that satisfies the proposition by taking a finite combination of the above cases; one corresponding to each type of intersection.

\end{proof}

\end{document}
